\documentclass[a4paper,12pt]{article}

\usepackage{iftex}
\ifPDFTeX
  \usepackage[T2A]{fontenc}
  \usepackage[utf8]{inputenc}
  \usepackage[russian]{babel}
  \usepackage{helvet}
  \usepackage{microtype}
\else
  \usepackage{fontspec}
  \setmainfont{DejaVu Sans}[AutoFakeSlant=0.18]
  \setsansfont{DejaVu Sans}[AutoFakeSlant=0.18]
\fi
\renewcommand{\familydefault}{\sfdefault}
\renewcommand{\today}{17 июля 2026 г.}
\usepackage{amsmath,amssymb,amsthm,mathtools}
\usepackage{icomma}
\usepackage[a4paper,left=18mm,right=18mm,top=18mm,bottom=20mm]{geometry}
\usepackage{tikz}
\usepackage{tkz-euclide}
\usepackage{pgfplots}
\pgfplotsset{compat=1.18}
\usepackage{tabularx,booktabs,longtable,array}
\usepackage{enumitem}
\usepackage{hyperref}
\usepackage{parskip}
\usepackage{fancyhdr}
\usepackage{setspace}
\usepackage{xcolor}

\definecolor{accent}{HTML}{7B3142}
\definecolor{darkaccent}{HTML}{4E2430}
\definecolor{textgray}{HTML}{30343B}
\definecolor{softgray}{HTML}{6E7279}

\hypersetup{colorlinks=true,linkcolor=darkaccent,urlcolor=accent}
\onehalfspacing
\setlength{\parindent}{0pt}
\setlength{\parskip}{5pt}
\setlength{\emergencystretch}{3em}
\setlist[itemize]{leftmargin=6mm,itemsep=2pt,topsep=3pt}
\setlist[enumerate]{leftmargin=8mm,itemsep=3pt,topsep=3pt}
\renewcommand{\arraystretch}{1.35}
\sloppy

\pagestyle{fancy}
\fancyhf{}
\lhead{\textcolor{textgray}{Алканы: строение и свойства}}
\rhead{\textcolor{textgray}{17.07.26}}
\cfoot{\textcolor{softgray}{\thepage}}
\setlength{\headheight}{15pt}

\newcommand{\key}[1]{\textcolor{accent}{\textbf{#1}}}
\newcommand{\chem}[1]{\ensuremath{\mathrm{#1}}}

\begin{document}

\begin{titlepage}
  \centering
  \vspace*{23mm}
  {\Large\textcolor{accent}{ЧЕК-ЛИСТ}\par}
  \vspace{9mm}
  {\fontsize{29}{35}\selectfont\bfseries Алканы:\\[2mm] строение, номенклатура\\[2mm] и химические свойства\par}
  \vspace{9mm}
  {\large Гибридизация, формулы, реакции замещения,\\дегидрирование и крекинг\par}
  \vfill
  {\large Лёвин Артём Александрович\\эксклюзивно для Ксении Клыковой\par}
  \vspace{8mm}
  {\large \today\par}
  \vfill
  \begin{tikzpicture}[remember picture,overlay]
    \draw[accent,line width=2pt]
      ([xshift=22mm,yshift=24mm]current page.south west)
      .. controls +(38mm,22mm) and +(-35mm,-10mm) ..
      ([xshift=-22mm,yshift=35mm]current page.south east);
    \draw[darkaccent,line width=.8pt]
      ([xshift=32mm,yshift=18mm]current page.south west)
      .. controls +(45mm,28mm) and +(-48mm,-14mm) ..
      ([xshift=-30mm,yshift=28mm]current page.south east);
  \end{tikzpicture}
\end{titlepage}

\section{Что нужно знать об алканах}

\key{Алканы} --- насыщенные ациклические углеводороды, в молекулах которых есть только одинарные связи \chem{C-C} и \chem{C-H}. Для них справедлива общая формула
\[
\boxed{\chem{C_{n}H_{2n+2}}},\qquad n\ge1.
\]

\begin{table}[ht]
\centering
\caption{Первые представители гомологического ряда}
\begin{tabularx}{\linewidth}{>{\bfseries}p{30mm} p{27mm} X}
\toprule
Название & Формула & Сокращённая структурная запись \\
\midrule
Метан & \chem{CH_4} & \chem{CH_4} \\
Этан & \chem{C_2H_6} & \chem{CH_3-CH_3} \\
Пропан & \chem{C_3H_8} & \chem{CH_3-CH_2-CH_3} \\
Бутан & \chem{C_4H_{10}} & \chem{CH_3-CH_2-CH_2-CH_3} \\
Пентан & \chem{C_5H_{12}} & \chem{CH_3-(CH_2)_3-CH_3} \\
Гексан & \chem{C_6H_{14}} & \chem{CH_3-(CH_2)_4-CH_3} \\
\bottomrule
\end{tabularx}
\end{table}

\key{Быстрая проверка формулы:} если в открытой цепи $n$ атомов углерода и только одинарные связи, число атомов водорода должно быть равно $2n+2$.

\section{Гибридизация и пространственное строение}

Каждый атом углерода в алкане образует четыре $\sigma$-связи и находится в состоянии $sp^3$-гибридизации. Четыре гибридные орбитали направлены к вершинам тетраэдра; идеальный валентный угол близок к $109{,}5^\circ$.

\begin{table}[ht]
\centering
\caption{Связь кратности и гибридизации}
\begin{tabularx}{\linewidth}{>{\bfseries}p{28mm} p{36mm} X}
\toprule
Тип & Признак & Геометрия вокруг атома углерода \\
\midrule
$sp^3$ & четыре $\sigma$-связи & тетраэдрическая \\
$sp^2$ & одна двойная связь & плоская тригональная \\
$sp$ & одна тройная или две двойные связи & линейная \\
\bottomrule
\end{tabularx}
\end{table}

В алканах встречается только $sp^3$-гибридизация углерода. Свободное вращение вокруг связи \chem{C-C} создаёт разные конформации, но не меняет порядок соединения атомов.

\subsection{Правило четырёхвалентности}

При построении структуры посчитайте связи у каждого атома \chem{C}. Одинарная связь считается за одну, двойная --- за две, тройная --- за три. Сумма должна быть равна четырём.

Например, в пентане
\[
\chem{CH_3-CH_2-CH_2-CH_2-CH_3}
\]
крайние атомы углерода несут по три атома водорода, а внутренние --- по два.

\section{Основы номенклатуры}

\begin{enumerate}
  \item Найдите самую длинную непрерывную углеродную цепь.
  \item Пронумеруйте её с того конца, к которому ближе первый заместитель.
  \item Назовите заместители с номерами их положений.
  \item Одинаковые заместители обозначьте приставками ди-, три-, тетра-.
  \item Завершите название именем главной цепи с суффиксом «-ан».
\end{enumerate}

Примеры:
\[
\chem{CH_3-CH_2-CH_2-CH_3}\quad\text{--- бутан},
\]
\[
\chem{CH_3-CH(CH_3)-CH_2-CH_3}\quad\text{--- 2-метилбутан}.
\]

\section{Горение}

При полном горении алкана образуются углекислый газ и вода:
\[
\chem{C_{n}H_{2n+2}}+\frac{3n+1}{2}\,\chem{O_2}
\longrightarrow n\,\chem{CO_2}+(n+1)\,\chem{H_2O}.
\]

Алгоритм уравнивания: сначала углерод, затем водород, последним кислород. Дробный коэффициент перед \chem{O_2} допустим как промежуточный; чтобы получить целые коэффициенты, умножьте всё уравнение на два.

\subsection{Пример: горение бутана}

\[
\chem{C_4H_{10}}+\frac{13}{2}\,\chem{O_2}\longrightarrow
4\,\chem{CO_2}+5\,\chem{H_2O},
\]
\[
\boxed{2\,\chem{C_4H_{10}}+13\,\chem{O_2}\longrightarrow
8\,\chem{CO_2}+10\,\chem{H_2O}}.
\]

\section{Реакции замещения}

\subsection{Галогенирование}

Под действием света или при нагревании атом водорода замещается галогеном:
\[
\chem{CH_4+Cl_2}\xrightarrow{h\nu}\chem{CH_3Cl+HCl}.
\]

Хлорирование метана может продолжаться ступенчато:
\[
\chem{CH_4}\longrightarrow\chem{CH_3Cl}\longrightarrow\chem{CH_2Cl_2}
\longrightarrow\chem{CHCl_3}\longrightarrow\chem{CCl_4}.
\]

В высших алканах обычно образуется смесь изомеров. Замещение преимущественно идёт у более замещённого атома углерода: третичный водород реакционноспособнее вторичного, а вторичный --- первичного. Особенно отчётливо эта селективность проявляется при бромировании; это не означает образования только одного продукта.

\subsection{Нитрование по Коновалову}

В школьной записи атом водорода замещается нитрогруппой:
\[
\chem{RH+HONO_2}\xrightarrow{t}\chem{RNO_2+H_2O}.
\]
Для метана:
\[
\chem{CH_4+HNO_3}\xrightarrow{t}\chem{CH_3NO_2+H_2O}.
\]

\subsection{Сульфирование: важное уточнение}

Обобщённую схему замещения записывают как
\[
\chem{RH+HOSO_3H}\longrightarrow\chem{RSO_3H+H_2O}.
\]
Однако низшие алканы химически малоактивны, и прямое сульфирование серной кислотой не является простой универсальной реакцией. Для высших алканов нужны жёсткие условия; в промышленности применяют также сульфохлорирование и сульфоокисление. Поэтому всегда ориентируйтесь на реагенты и условия конкретной задачи.

\section{Реакции с изменением углеродного скелета или кратности}

\subsection{Дегидрирование}

При нагревании над катализатором от алкана отщепляется водород и образуется алкен:
\[
\chem{C_4H_{10}}\xrightarrow{\text{кат.},\,t}\chem{C_4H_8+H_2}\!\uparrow.
\]
Для бутана возможны бут-1-ен и бут-2-ен. Положение двойной связи определяют условия и формулировка задачи.

\subsection{Крекинг}

При сильном нагревании без доступа кислорода длинная цепь расщепляется. Один из возможных вариантов:
\[
\chem{C_5H_{12}}\xrightarrow{t}\chem{C_2H_4+C_3H_8}.
\]
Суммарное число атомов каждого элемента сохраняется. В типичной школьной схеме продуктами являются более короткий алкан и алкен. После записи обязательно проверьте четыре связи у каждого атома углерода.

\subsection{Термическое разложение метана}

\[
\chem{CH_4}\xrightarrow{1000\,{}^\circ\mathrm{C}}\chem{C+2H_2}\!\uparrow,
\]
\[
2\,\chem{CH_4}\xrightarrow{1500\,{}^\circ\mathrm{C}}\chem{C_2H_2+3H_2}\!\uparrow.
\]

\section{Типичные ошибки}

\begin{enumerate}[itemsep=1pt]
  \item Использовать формулу \chem{C_{n}H_{2n+2}} для циклоалкана или алкена.
  \item Оставлять у атома углерода три или пять связей.
  \item Забывать свет или нагревание над стрелкой галогенирования.
  \item Считать преимущественный продукт единственно возможным.
  \item Уравнивать кислород до углерода и водорода в реакции горения.
  \item При крекинге терять атомы водорода или нарушать баланс углерода.
  \item Писать ацетилен из одной молекулы метана: для \chem{C_2H_2} нужны две молекулы \chem{CH_4}.
\end{enumerate}

\section{Итоговый чек-лист}

\begin{itemize}[label=$\square$,itemsep=1pt]
  \item Я узнаю алкан по одинарным связям и формуле \chem{C_{n}H_{2n+2}}.
  \item Я определяю $sp^3$-гибридизацию и тетраэдрическую геометрию.
  \item Я проверяю четыре связи у каждого атома углерода.
  \item Я называю главную цепь и нумерую заместители.
  \item Я уравниваю горение в порядке \chem{C}\,$\rightarrow$\,\chem{H}\,$\rightarrow$\,\chem{O}.
  \item Я указываю условия галогенирования, дегидрирования и крекинга.
  \item Я различаю замещение, отщепление и расщепление цепи.
  \item Я проверяю баланс атомов во всех продуктах.
\end{itemize}

\newpage
\section{Тренировочный блок}

\textbf{Средний уровень}
\begin{enumerate}
  \item Запишите молекулярную формулу алкана с семью атомами углерода.
  \item Определите формулу алкана, если в его молекуле 22 атома водорода.
  \item Назовите соединение \chem{CH_3-CH_2-CH_2-CH_3}.
\end{enumerate}

\textbf{Уровень выше среднего}
\begin{enumerate}[resume]
  \item Уравняйте полное горение пропана.
  \item Запишите монохлорирование этана и укажите условие.
  \item Запишите нитрование метана по Коновалову.
  \item Запишите общую схему дегидрирования бутана.
  \item Дополните схему крекинга: \chem{C_6H_{14}}\,$\rightarrow$\,\chem{C_2H_4}+\ldots
  \item Для 2-метилбутана запишите преимущественный продукт монохлорирования у третичного атома углерода; помните, что реально образуется смесь.
\end{enumerate}

\textbf{Сложный уровень}
\begin{enumerate}[resume]
  \item При полном горении одного моля алкана образовалось 6 моль \chem{CO_2} и 7 моль \chem{H_2O}. Определите алкан и запишите уравнение с целыми коэффициентами.
\end{enumerate}

\newpage
\section{Ответы}

\begin{table}[ht]
\centering
\caption{Ответы к тренировочному блоку}
\begin{tabularx}{\linewidth}{>{\bfseries}p{11mm} X}
\toprule
№ & Ответ \\
\midrule
1 & \chem{C_7H_{16}} \\
2 & $2n+2=22$, $n=10$; \chem{C_{10}H_{22}} \\
3 & Бутан \\
4 & \chem{C_3H_8+5O_2}\,$\rightarrow$\,\chem{3CO_2+4H_2O} \\
5 & \chem{C_2H_6+Cl_2}\,$\xrightarrow{h\nu}$\,\chem{C_2H_5Cl+HCl} \\
6 & \chem{CH_4+HNO_3}\,$\xrightarrow{t}$\,\chem{CH_3NO_2+H_2O} \\
7 & \chem{C_4H_{10}}\,$\xrightarrow{\text{кат.},\,t}$\,\chem{C_4H_8+H_2} \\
8 & \chem{C_4H_{10}}; полная схема: \chem{C_6H_{14}}\,$\rightarrow$\,\chem{C_2H_4+C_4H_{10}} \\
9 & \chem{CH_3-C(Cl)(CH_3)-CH_2-CH_3} --- 2-хлор-2-метилбутан \\
10 & Гексан, \chem{C_6H_{14}}; \chem{2C_6H_{14}+19O_2}\,$\rightarrow$\,\chem{12CO_2+14H_2O} \\
\bottomrule
\end{tabularx}
\end{table}

\section*{Домашний минимум}
\begin{enumerate}
  \item \chem{CH_4+Cl_2}\,$\xrightarrow{h\nu}$\,\chem{CH_3Cl+HCl}.
  \item \chem{CH_4+HNO_3}\,$\xrightarrow{t}$\,\chem{CH_3NO_2+H_2O}.
  \item \chem{2CH_4}\,$\xrightarrow{1500\,{}^\circ\mathrm{C}}$\,\chem{C_2H_2+3H_2}.
\end{enumerate}

\vfill
\begin{center}
\textcolor{softgray}{Каждую реакцию завершайте двумя проверками: баланс атомов и четыре связи у углерода.}
\end{center}

\end{document}
